
A finalidade deste capítulo é delimitar o que se pretende construir ao longo das duas fases do Trabalho de Conclusão de Curso e, em seguida, apresentar a cadeia temporal de atividades prevista na pré-proposta aprovada. A descrição parte da plataforma vista como produto educacional, atravessa as macroatividades técnicas necessárias à sua materialização e culmina na exposição do cronograma e do regime de trabalho em dupla adotado.

O artefato que orienta este trabalho não é apenas um compilador ou um editor de código, mas sim uma plataforma educacional web cuja unidade de valor é a possibilidade de o estudante personalizar a própria linguagem que utilizará para aprender programação. Em termos concretos, oferece-se ao usuário a possibilidade de redefinir o vocabulário das palavras-chave, escolher os delimitadores de bloco que melhor representem sua noção de escopo, substituir operadores lógicos e relacionais por palavras de uso natural, reescrever os literais booleanos no idioma de sua preferência e ajustar as regras de finalização de instrução. Cada uma dessas dimensões responde, em alguma medida, a uma fonte conhecida de atrito sintático no contato inicial com linguagens consolidadas no mercado, sendo a sua reunião sob uma mesma interface o que sustenta o caráter inovador da proposta.

Conforme definido na pré-proposta, a metodologia foi organizada em seis macroetapas complementares. Essas etapas não devem ser entendidas como blocos completamente isolados, pois o desenvolvimento de uma plataforma interativa exige revisões sucessivas entre interface, interpretador e experiência de uso. Ainda assim, elas delimitam com clareza o percurso técnico e acadêmico do projeto:

\begin{itemize}
    \item \textbf{Revisão bibliográfica:} levantamento das bases conceituais sobre interpretadores, compiladores e ferramentas educacionais relacionadas, com o objetivo de sustentar tanto as decisões de projeto quanto a posterior análise comparativa;
    \item \textbf{Definição de requisitos e arquitetura:} consolidação dos elementos configuráveis da linguagem, da modelagem das camadas da plataforma e dos contratos de comunicação entre o núcleo do interpretador e a IDE;
    \item \textbf{Desenvolvimento da plataforma:} implementação da aplicação web, contemplando os fluxos de edição, navegação, personalização, autenticação e uso educacional da IDE;
    \item \textbf{Implementação da análise dinâmica léxica e sintática:} integração entre o código escrito na IDE e as fases de reconhecimento de \textit{tokens} e validação sintática com suporte às regras configuradas pelo usuário;
    \item \textbf{Geração e interpretação do código intermediário:} disponibilização, no ambiente da IDE, da execução do código intermediário e do retorno de avisos, erros e resultados ao usuário;
    \item \textbf{Integração e testes entre plataforma e interpretador:} validação dos fluxos principais, com atenção à consistência da experiência, à acessibilidade, à responsividade e à preparação para os estudos com usuários.
\end{itemize}

A distribuição temporal dessas atividades ao longo dos doze meses que compreendem o TCC~I e o TCC~II é apresentada na Tabela~\ref{tab:cronograma}, conforme o cronograma registrado na pré-proposta. Cada coluna numerada corresponde a um mês corrido das duas etapas e cada célula marcada com ``X'' sinaliza a previsão de execução da respectiva atividade no período indicado.

\begin{table}[H]
    \centering
    \caption{Cronograma de desenvolvimento do trabalho de conclusão de curso.}
    \label{tab:cronograma}
    \footnotesize
    \setlength{\tabcolsep}{4pt}
    \renewcommand{\arraystretch}{1.2}
    \resizebox{\textwidth}{!}{%
    \begin{tabular}{|p{6.5cm}|c|c|c|c|c|c|c|c|c|c|c|c|}
        \hline
        \multirow{2}{*}{\textbf{Atividade}} & \multicolumn{6}{c|}{\textbf{TCC I}} & \multicolumn{6}{c|}{\textbf{TCC II}} \\
        \cline{2-13}
         & \textbf{01} & \textbf{02} & \textbf{03} & \textbf{04} & \textbf{05} & \textbf{06} & \textbf{07} & \textbf{08} & \textbf{09} & \textbf{10} & \textbf{11} & \textbf{12} \\
        \hline
        Revisão Bibliográfica & X & X & X & X & & & & & & & & \\
        \hline
        Escrita da Introdução e Revisão de Literatura & & X & X & X & & & & & & & & \\
        \hline
        Predefinição de escopo e Especificação dos comandos dinâmicos & & & X & X & & & & & & & & \\
        \hline
        Escrita da Metodologia & & & X & & X & X & & & & & & \\
        \hline
        Arquitetura e Modelagem da Plataforma & & & X & X & & & & & & & & \\
        \hline
        Customização dinâmica de tokens e sintaxe & & & & X & X & & & & & & & \\
        \hline
        Execução da Análise Léxica e Sintática com o código da IDE & & & & & X & X & & & & & & \\
        \hline
        Execução do código intermediário no Terminal da IDE com \textit{feedback} (avisos e erros) & & & & & & X & X & X & & & & \\
        \hline
        Finalização de fluxos, acessibilidade e responsividade & & & & & & X & X & X & & & & \\
        \hline
        Testes, Validação com Usuários e Coleta de \textit{Feedback} & & & & & & & & & X & X & X & \\
        \hline
        Escrita do TCC (Resultados, Discussão e Conclusão) & & & & & & & & & X & X & X & X \\
        \hline
    \end{tabular}%
    }
    \legend{Fonte: elaborado pelo autor.}
\end{table}

Convém esclarecer, neste ponto, o regime de trabalho adotado, dada sua influência direta sobre o escopo do projeto e sobre os capítulos subsequentes. A pré-proposta aprovada para este Trabalho de Conclusão de Curso prevê a \textbf{participação conjunta em todas as etapas} do desenvolvimento, mas com funções complementares. Em respeito a esse compromisso, as decisões de arquitetura, os ciclos de revisão de código e as etapas de integração foram conduzidos de modo colaborativo pelos dois integrantes da dupla.

A partir dessa base comum, a divisão complementar de eixos de aprofundamento manifesta-se nos capítulos de redação acadêmica e nas evoluções incrementais subsequentes. No presente documento, o autor concentra a investigação na plataforma e na customização de comandos: arquitetura da aplicação web, concepção do assistente de personalização, experiência de uso da IDE, execução do código no terminal integrado e procedimentos de coleta de \textit{feedback} junto aos estudantes. O detalhamento aprofundado da evolução do núcleo do interpretador, incluindo a geração de código intermediário e sua integração técnica ao restante da plataforma, é desenvolvido em paralelo pelo discente Victor de Souza Vilela da Silva. Essa divisão de eixos não fragmenta o produto: o que se distribui é a profundidade analítica de cada autor sobre partes específicas do mesmo artefato, conduzido pelos dois em todas as etapas. Tal arranjo justifica-se pela amplitude do escopo, que, individualmente, tenderia a resultar em um interpretador simplificado ou em uma plataforma incompleta, restrita a um terminal.

Concluída a etapa coberta por este documento, dedicada à fundamentação teórica, à modelagem da plataforma, à customização dinâmica de \textit{tokens} e sintaxe, à execução das análises léxica e sintática a partir da IDE e à execução do código intermediário no terminal integrado, a fase subsequente direcionará seus esforços à condução dos estudos de caso, à coleta de evidências empíricas em sala de aula e à análise comparativa entre a plataforma desenvolvida e as ferramentas discutidas no referencial teórico. O próximo capítulo, dedicado à modelagem, descreve o estado em que o sistema se encontra ao término do TCC~I.
